\documentclass{article}
\usepackage[utf8]{inputenc}
\usepackage[a4paper, margin=1in]{geometry}
\usepackage{hyperref}
\usepackage{amsmath, amssymb, amsthm}
\usepackage{enumitem}

\hypersetup{
  pdftitle={Self-Check Questions},
  pdfauthor={Mathematics Study Notes},
  pdfsubject={Mathematics},
  pdfkeywords={Mathematics, Self-Check Questions, Active Recall},
  pdfcreator={LaTeX},
  pdfproducer={pdfLaTeX}
}

\hypersetup{
  colorlinks=true,
  linkcolor=blue,
  urlcolor=blue
}

\urlstyle{same}

\title{Self-Check Questions}
\author{Mathematics Study Notes}
\date{\today}

\setcounter{tocdepth}{1}

\begin{document}

\maketitle

\section{Introduction}

Mathematics is not a subject that can be mastered through passive reading alone. True understanding requires the ability to retrieve concepts, definitions, and problem-solving strategies from memory. Self-check questions provide structure for this: a curated list of prompts that guide the learner through the essential components of a topic.

\subsection{What Are Self-Check Questions?}

A self-check question is a prompt that the learner attempts to answer, derive, or solve from memory, without looking at the solution. Instead of rereading notes or textbooks, the learner tries to reconstruct the material directly.

In mathematics, this might include:
\begin{itemize}
  \item Stating definitions (e.g., convergence of a series),
  \item Recalling theorems (e.g., the Integral Test),
  \item Re-deriving formulas,
  \item Solving representative problems from memory.
\end{itemize}

\subsection{Role of Self-Check Questions}

For example, in studying infinite series, one might ask:
\begin{itemize}
  \item What is the definition of convergence?
  \item When can the Integral Test be applied?
  \item Why does removing finitely many terms not affect convergence?
\end{itemize}

Such questions ensure that the learner engages with both conceptual understanding and procedural knowledge.

\subsection{Why It Works}

Answering self-check questions relies on \emph{active recall}: deliberately retrieving information from memory rather than re-exposing oneself to it. Active recall strengthens memory through effortful retrieval. Each time information is successfully recalled, the neural pathways associated with that knowledge become more robust. This leads to:
\begin{itemize}
  \item Improved long-term retention,
  \item Faster recall during problem solving,
  \item Better ability to connect different concepts.
\end{itemize}

Moreover, attempting to answer questions reveals gaps in understanding. These gaps are often invisible during passive review but become immediately apparent during recall attempts.

\clearpage

\tableofcontents

\clearpage

\input{openstax/_index}
\input{mc/_index}
\input{misc/_index}
\input{numeric/_index}

\clearpage

\input{self-check-stats}

\end{document}
